\section{矩阵乘法补充：Tiling 与 L2 Cache 优化}

\subsection{为什么手写矩阵乘法内核？}

PyTorch 内置的矩阵乘法（调用 cuBLAS/CUTLASS）已经高度优化。那为什么还要手写？

\textbf{答案：融合}。当你需要计算 $\text{GeLU}(A \cdot B)$ 时，朴素做法是先算 $A \cdot B$（写回全局内存），再算 GeLU（从全局内存读取）。如果你有自己的矩阵乘法内核，就可以在同一个内核中完成矩阵乘法\textbf{和} GeLU，避免中间结果写回全局内存。

\subsection{Tiling 策略}

矩阵乘法的 Tiling 策略（详见 Lecture 5）：
\begin{enumerate}
    \item 将 $A$、$B$ 矩阵切分为小块（tile）
    \item 将 tile 加载到 Shared Memory
    \item 在 Shared Memory 中执行小矩阵乘法
    \item 累加部分和
    \item 写回结果
\end{enumerate}

\subsection{L2 Cache 友好的遍历顺序}

\begin{figure}[H]
\centering
\begin{tikzpicture}[scale=0.45]
% Row-major order
\node[font=\bfseries\small] at (1.5, 4.5) {行优先遍历};
\foreach \i in {0,...,2} {
    \foreach \j in {0,...,2} {
        \draw (\j, 2-\i) rectangle (\j+1, 3-\i);
    }
}
\draw[->, red, thick] (0.5, 2.5) -- (1.5, 2.5) -- (2.5, 2.5);
\draw[->, red, thick] (0.5, 1.5) -- (1.5, 1.5) -- (2.5, 1.5);
\draw[->, red, thick] (0.5, 0.5) -- (1.5, 0.5) -- (2.5, 0.5);
\node[font=\scriptsize] at (1.5, -0.3) {加载 9+81=90 个块};

% Grouped order
\node[font=\bfseries\small] at (7.5, 4.5) {分组遍历};
\foreach \i in {0,...,2} {
    \foreach \j in {0,...,2} {
        \draw (6+\j, 2-\i) rectangle (6+\j+1, 3-\i);
    }
}
\draw[->, green!60!black, thick] (6.5, 2.5) -- (6.5, 1.5) -- (6.5, 0.5);
\draw[->, green!60!black, thick] (7.5, 2.5) -- (7.5, 1.5) -- (7.5, 0.5);
\draw[->, green!60!black, thick] (8.5, 2.5) -- (8.5, 1.5) -- (8.5, 0.5);
\node[font=\scriptsize] at (7.5, -0.3) {加载 27+27=54 个块};
\end{tikzpicture}
\caption{不同遍历顺序对 L2 Cache 命中率的影响}
\end{figure}

分组遍历（grouped ordering）让相邻的 tile 共享更多的 $A$ 或 $B$ 矩阵块，从而提高 L2 Cache 命中率。

\subsection{本章小结}

矩阵乘法是 GPU 上最受优化的操作，但手写内核仍有价值——主要用于与其他操作的融合。Tiling 和 L2 Cache 友好的遍历顺序是两个关键优化策略。

%% ========== Section 10.5: 实践建议 ==========

